\scanpage{91}
% archive foliation: 92
4. \deleted{\unsure{Théorè}} \uline{Déformations de groupes de} BT : \uline{énoncés}.

Le résultat principal, qui sera essentiel pour la théorie de Dieudonné est le suivant(cf lettre Illusie, 6.4):

\uline{Théorème} 4.1. Soit $S_o\hookrightarrow S$ une nilimmersion, avec $S$ affine, $G_o$ un groupe de BT sur $S_o$. Alors:

a) Il existe un groupe de BT $G$ sur $S$ qui prolonge $G_o$.

b) Lapplication
\[
(*)\qquad E(G_o,S)\longrightarrow E(G_o(n),S)
\]
est surjective\unsure{;|.}

c) Si $i$ est nilpotente d'ordre $k$ et si $p^N1_{S_o}=0$, alors l'application $(*)$ est bijective pour $n\geq kN$.

d) Si de plus $k=1$, alors $E(G_o,S)$ est un torseur sous $\underline t^{\deleted{\unsure{*}}}_{G_o^*}\otimes t_{G_o}$.

Ce résultat s'obtient via un résultat de déformation pour les groupes de BT tronqués (je ne sais pas procéder directement !).

5. \uline{Rapport sur le complexe cotangent relatif}.
